\section{Main theorem}

The core result is finite-stage boundary localization. It does not require a likelihood-ratio bound that is uniform over all possible clinical-history lengths, nor does it require an unbounded clinical-evidence tail. At any fixed observed history length, the relevant quantity is simply the high-probability scale of the residual genomic log-likelihood ratio. Stronger sequential and extreme-tail statements are separated below as conditional extensions.

\paragraph*{Notation.} For sequential clinical histories $C_n=(X_1,\ldots,X_n)$, write
\begin{equation}
R_{G,n}=\log \LR_{G\mid C_n}(G;C_n)
\end{equation}
for the residual genomic log-likelihood ratio at stage $n$. For a clinical likelihood-ratio threshold $\lambda>0$, define
\begin{equation}
D_{C_n}=\ind\{\LR_{C_n}\ge \lambda\},\qquad
D_{C_n,G}=\ind\{\LR_{C_n}\LR_{G\mid C_n}\ge \lambda\}.
\end{equation}
For each fixed $n$ and tolerance $\delta\in(0,1)$, define the finite-stage residual-evidence scale
\begin{equation}
b_{\delta,n}
=\inf\left\{b\ge0:\Pzero\bigl(|R_{G,n}|>b\bigr)\le\delta\right\},
\label{eq:finite-stage-bound}
\end{equation}
whenever this quantile is finite. Thus $b_{\delta,n}$ is a property of the residual genomic evidence distribution at the clinical stage actually being analyzed; no uniformity over $n$ is required.

\begin{theorem}[Finite-stage boundary localization]
\label{thm:main}
Assume the relevant likelihood ratios are well defined and $b_{\delta,n}<\infty$ at the finite stage under consideration.

\emph{(i) Exact factorization.} For every $n$,
\begin{equation}
\LR_{C_n,G}(c_n,g)=\LR_{C_n}(c_n)\LR_{G\mid C_n}(g;c_n).
\end{equation}

\emph{(ii) High-probability boundary localization.} For every $n$, every decision threshold $\lambda>0$, and every $\delta\in(0,1)$,
\begin{equation}
\boxed{
\Pzero\Bigl(
D_{C_n,G}\neq D_{C_n},\
\bigl|\log \LR_{C_n}-\log\lambda\bigr|>b_{\delta,n}
\Bigr)\le \delta.
}
\label{eq:stochastic-band-result}
\end{equation}
Thus, at a fixed clinical stage, genomic information can reverse the clinical decision outside a log-LR neighborhood of half-width $b_{\delta,n}$ only on a control-side exceptional event of probability at most $\delta$.
\end{theorem}

\begin{proof}
Part (i) is the chain rule:
\begin{align}
\LR_{C_n,G}(c_n,g)
&=\frac{p(c_n,g\mid Y=1)}{p(c_n,g\mid Y=0)}\\
&=\frac{p(c_n\mid Y=1)p(g\mid Y=1,c_n)}{p(c_n\mid Y=0)p(g\mid Y=0,c_n)}\\
&=\LR_{C_n}(c_n)\LR_{G\mid C_n}(g;c_n).
\end{align}
For part (ii), if both $|\log \LR_{C_n}-\log\lambda|>b_{\delta,n}$ and $|R_{G,n}|\le b_{\delta,n}$ hold, then adding $R_{G,n}$ cannot change the sign of $\log \LR_{C_n}-\log\lambda$. Hence
\begin{multline}
\Bigl\{D_{C_n,G}\neq D_{C_n},\
|\log \LR_{C_n}-\log\lambda|>b_{\delta,n}\Bigr\}\\
\subseteq \{|R_{G,n}|>b_{\delta,n}\},
\end{multline}
whose probability is at most $\delta$ by the definition of $b_{\delta,n}$.
\end{proof}

\subsection{Stronger finite-stage special cases}

\begin{corollary}[Deterministic bounded-residual special case]
\label{cor:deterministic-band}
Fix a clinical representation $C$. If there exist constants $0<m\le1\le M<\infty$ such that
\begin{equation}
0<m\le \LR_{G\mid C}(g;c)\le M<\infty
\label{eq:boundedconditional}
\end{equation}
for $\Pzero$-almost every relevant $(g,c)$, then genomic information can change which side of threshold $\lambda$ an observation occupies only if
\begin{equation}
\boxed{\frac{\lambda}{M}\le \LR_C(c)\le \frac{\lambda}{m}.}
\label{eq:boundaryband}
\end{equation}
Equivalently,
\begin{equation}
\boxed{-\log M\le \log \LR_C(c)-\log\lambda\le -\log m.}
\label{eq:logboundaryband}
\end{equation}
The deterministic decision-relevant band has fixed width $\log(M/m)$ on the clinical log-likelihood scale. If the same $m$ and $M$ apply uniformly to every $C_n$, then the optional uniform-tightness condition introduced below holds with $b_\delta=\max\{\log M,-\log m\}$.
\end{corollary}
The proof is given in Appendix~\ref{app:proofs}.

\begin{corollary}[Posterior-risk form for a calibrated clinical score]
\label{cor:zebra}
If a clinical score $Z$ is calibrated so that $Z=\mathbb{P}(Y=1\mid C)$, let $\tau\in(0,1)$ be any decision threshold and define posterior-odds threshold $\kappa=\tau/(1-\tau)$. Under Corollary~\ref{cor:deterministic-band}, genomic information can change the decision only if
\begin{equation}
\boxed{\frac{\kappa}{M}\le \frac{Z}{1-Z}\le \frac{\kappa}{m}.}
\label{eq:zboundary}
\end{equation}
Equivalently,
\begin{equation}
-\log M\le \operatorname{logit}(Z)-\operatorname{logit}(\tau)\le -\log m.
\end{equation}
This is an idealized calibrated-score corollary, not the empirical representation used in the Colorado analysis. The raw ZeBRA output is not calibrated as a posterior probability there; the empirical analysis instead estimates a score-level likelihood ratio directly.
\end{corollary}
The proof is given in Appendix~\ref{app:proofs}.

\subsection{Conditional sequential and extreme-tail extensions}

The following assumptions are \emph{not} required for Theorem~\ref{thm:main}. They are introduced only for stronger statements about what happens as clinical history lengthens or at arbitrarily extreme clinical likelihood-ratio thresholds.

\textbf{A1-U (Optional uniform tightness of residual genomic evidence).} For each tolerance $\delta\in(0,1)$, suppose there exists a finite $b_\delta$, independent of $n$, such that
\begin{equation}
\sup_n\Pzero\bigl(|R_{G,n}|>b_\delta\bigr)\le\delta.
\label{eq:stochasticbound}
\end{equation}
This is substantially stronger than finite-stage boundedness: it requires one residual-genomic scale to remain valid uniformly as the clinical record lengthens. The Colorado analysis does not test this condition.

\textbf{A2 (Unbounded positive clinical-evidence tail).} Suppose that for arbitrarily large $K>0$ there exists an $n$ such that
\begin{equation}
\Pzero\bigl(\log \LR_{C_n}(C_n)\ge K\bigr)>0.
\label{eq:clinicaltailassumption}
\end{equation}
A2 is an explicit structural assumption that arbitrarily strong positive clinical evidence is attainable on rare trajectories. Sequential factorization can explain how evidence accumulates, but it does not imply A2. Empirical data over a finite range can be consistent with this geometry but cannot prove mathematical unboundedness.

\begin{center}
\footnotesize
\begin{tabular}{@{}lp{0.68\columnwidth}@{}}
\toprule
Notation & Meaning \\
\midrule
$b_{\delta,n}$ & Finite-stage $(1-\delta)$ residual-genomic log-LR scale at a particular clinical-history stage $n$; used in the main theorem. \\
$b_{\delta}$ & A single $n$-independent scale available only under the stronger A1-U condition. \\
\bottomrule
\end{tabular}
\end{center}

\begin{theorem}[Conditional sequential clinical-tail dominance]
\label{thm:tail-extension}
Assume A1-U. Then for every $n$, every positive upper-tail threshold $K>0$, and every $\delta\in(0,1)$,
\begin{equation}
\boxed{
\Pzero\left(
\log\LR_{C_n}\ge K,\
\left|
\frac{\log\LR_{C_n,G}-\log\LR_{C_n}}
     {\log\LR_{C_n}}
\right|>\frac{b_\delta}{K}
\right)\le\delta.
}
\label{eq:taildom-prob}
\end{equation}
If A2 also holds, positive-probability upper clinical-tail events exist for arbitrarily large $K$, while $b_\delta/K\to0$. Thus, conditional on A1-U and A2, residual genomic log-evidence becomes relatively negligible on sufficiently extreme positive clinical-evidence tails except on a control-side event of probability at most $\delta$.
\end{theorem}
The proof is given in Appendix~\ref{app:proofs}. This theorem is a conditional sequential extension of Theorem~\ref{thm:main}, not an assumption needed for finite-stage localization.

\begin{proposition}[Genomic-only power bound and conditional extreme-specificity clinical superiority]
\label{prop:genomicpower}
Assume the marginal genomic likelihood ratio
\begin{equation}
\LR_G(g)=\frac{p(g\mid Y=1)}{p(g\mid Y=0)}
\end{equation}
is bounded above by $M_G<\infty$. Then:

\emph{(i)} for any genomic-only test $\phi(G)\in[0,1]$ with false-positive rate $\alpha=\Ezero[\phi(G)]$,
\begin{equation}
\boxed{\Eone[\phi(G)]\le M_G\alpha.}
\label{eq:genomicpowerbound}
\end{equation}

\emph{(ii)} if A2 holds, then for any $K>M_G$ with $\alpha_K=\Pzero(\LR_C\ge K)>0$, the clinical likelihood-ratio rule
\begin{equation}
\phi_K(C)=\ind\{\LR_C(C)\ge K\}
\end{equation}
has
\begin{equation}
\frac{\Pone(\LR_C\ge K)}{\Pzero(\LR_C\ge K)}\ge K>M_G,
\label{eq:clinicalLRplus}
\end{equation}
so at the same FPR $\alpha_K$ its power exceeds that of every genomic-only rule. Moreover $\alpha_K\le 1/K\to0$ as $K\to\infty$.
\end{proposition}
The proof is given in Appendix~\ref{app:proofs}. Part (i) is finite and does not require A2; only the arbitrarily extreme comparison in part (ii) invokes the unbounded-tail condition.

\begin{corollary}[Genomic post-processing inherits the marginal bound]
\label{cor:genomic-postprocessing}
Let $S$ be generated from the measured genomic state $G$ by a class-independent Markov kernel $Q(ds\mid g)$, so that $Y\to G\to S$ forms a Markov chain. This includes deterministic scores $S=f(G)$ and randomized post-processing whose additional randomness is independent of disease class conditional on $G$. If $\LR_G\le M_G$ almost surely under $\Pzero$, then $\LR_S\le M_G$ almost surely. Thus no such post-processing of the same genomic information can evade Proposition~\ref{prop:genomicpower}.
\end{corollary}
The proof is given in Appendix~\ref{app:proofs}.

\subsection{Local decision gain and global discrimination}

\begin{proposition}[Local decision gain can coexist with negligible global AUC change]
\label{prop:auc-local}
Let $S(C)$ be a clinical score and $S'(C,G)$ an augmented score. Suppose $S'=S$ except on a measurable region $B$ of the clinical--genomic state space. Then
\begin{align}
\left|\operatorname{AUC}(S')-\operatorname{AUC}(S)\right|
&\le \Pone(B)+\Pzero(B)-\Pone(B)\Pzero(B)\\
&\le \Pone(B)+\Pzero(B).
\label{eq:auc-local-bound}
\end{align}
Hence an augmented rule can produce a substantial change in decisions within a small boundary region while producing an arbitrarily small change in global AUC as the case and control mass of that region tends to zero.
\end{proposition}
The proof is given in Appendix~\ref{app:proofs}.
